\chapter{Postać Jordana}
\label{ch:lin17-jordan}
\index{postać Jordana@Postać Jordana}
\index{blok Jordana@Blok Jordana}
\index{łańcuch Jordana@Łańcuch Jordana}
\index{operator nilpotentny@Operator nilpotentny}

\noindent\textbf{Cel rozdziału.}
W Rozdziale~\ref{ch:lin16-niezmiennicze}
wyprowadziliśmy rozkład na uogólnione przestrzenie własne.
Teraz opiszemy dokładniej działanie operatora
na każdej z tych przestrzeni.
Dla endomorfizmu o rozszczepionym wielomianie
charakterystycznym uzyskamy postać Jordana.
Udowodnimy istnienie tej postaci, określimy
jej jednoznaczność z dokładnością do kolejności bloków
i pokażemy, jak odczytywać rozmiary bloków
z wymiarów jąder kolejnych potęg.

W całym rozdziale \(K\) oznacza dowolne ciało,
\(V\) skończenie wymiarową przestrzeń nad \(K\),
a \(T\in\operatorname{End}_K(V)\).
Nie zakładamy istnienia iloczynu skalarnego.
Przyjmujemy też przypadek \(V=\{0\}\),
w którym stosujemy bazę pustą i macierz pustą.

\section{Nilpotentność i łańcuchy Jordana}
\label{sec:lin17-nilpotent}

\begin{definicja}[Operator nilpotentny i indeks nilpotentności]
\label{def:lin17-nilpotent}
Endomorfizm \(N\in\operatorname{End}_K(V)\)
nazywamy \emph{nilpotentnym}, jeżeli
istnieje liczba całkowita \(m\ge1\), dla której
\(N^m=0\).
Jeżeli \(V\ne\{0\}\), najmniejszą taką
liczbę \(m\) nazywamy \emph{indeksem nilpotentności}.
Dla \(V=\{0\}\) przyjmujemy indeks \(0\).
\end{definicja}

\begin{stwierdzenie}[Jądra kolejnych potęg]
\label{prop:lin17-kernel-filtration}
Niech \(N\in\operatorname{End}_K(V)\)
będzie nilpotentny i niech \(n=\dim V\).
Dla \(k\ge0\) oznaczmy
\(K_k=\ker N^k\), przy czym \(N^0=I_V\).
Wówczas
\[
 \{0\}=K_0\subseteq K_1\subseteq
 K_2\subseteq\cdots\subseteq V.
\]
Jeżeli \(K_{k+1}=K_k\) dla pewnego \(k\ge0\),
to \(K_{k+j}=K_k\) dla każdego \(j\ge0\).
W szczególności dla \(V\ne\{0\}\)
indeks nilpotentności jest nie większy niż \(n\).
\end{stwierdzenie}
\begin{proof}
Z równości \(N^kv=0\) wynika
\(N^{k+1}v=N(N^kv)=0\), zatem
\(K_k\subseteq K_{k+1}\).
Jeżeli \(K_{k+1}=K_k\)
i \(N^{k+2}v=0\), to
\(Nv\in K_{k+1}=K_k\), więc
\(v\in K_{k+1}=K_k\).
Wobec tego \(K_{k+2}=K_{k+1}\).
Powtarzając argument, otrzymujemy
stabilizację wszystkich dalszych jąder.

Ponieważ \(N\) jest nilpotentny,
dla pewnego \(m\) mamy \(K_m=V\).
Jeśli \(V\ne\{0\}\), to
\(K_0=\{0\}\ne V\), więc
przed osiągnięciem \(V\)
nie może wystąpić równość kolejnych jąder.
Zatem w ciągu
\(K_0\subsetneq K_1\subsetneq\cdots\subsetneq V\)
każdy krok zwiększa wymiar przynajmniej o jeden.
Liczba ścisłych kroków nie przekracza \(n\),
czyli \(N^n=0\).
\end{proof}

\begin{definicja}[Łańcuch Jordana operatora nilpotentnego]
\label{def:lin17-jordan-chain}
Niech \(N\) będzie nilpotentnym
endomorfizmem \(V\).
\emph{Łańcuchem Jordana długości \(s\ge1\)}
nazywamy uporządkowaną rodzinę
\((v_1,\ldots,v_s)\) taką, że
\[
 v_1\ne0,\qquad Nv_1=0,\qquad
 Nv_j=v_{j-1}\quad(2\le j\le s).
\]
Wektor \(v_s\) nazywamy
\emph{wektorem końcowym} łańcucha.
\end{definicja}

\begin{stwierdzenie}[Niezależność wektorów łańcucha]
\label{prop:lin17-chain-independent}
Każdy łańcuch Jordana \((v_1,\ldots,v_s)\)
składa się z liniowo niezależnych wektorów,
a jego rozpiętość jest niezmiennicza względem \(N\).
W tej bazie macierz ograniczenia \(N\)
ma zera poza pierwszą nadprzekątną,
na której znajdują się jedynki.
\end{stwierdzenie}
\begin{proof}
Niech \(\sum_{j=1}^{s}a_jv_j=0\).
Zastosowanie \(N^{s-1}\) daje
\(a_sv_1=0\), a więc \(a_s=0\).
Następnie zastosowanie \(N^{s-2}\)
do pozostałej równości daje
\(a_{s-1}v_1=0\).
Kontynuując w ten sposób, dostajemy
\(a_j=0\) dla każdego \(j\).
Niezmienniczość wynika bezpośrednio
z warunków \(Nv_1=0\) i \(Nv_j=v_{j-1}\).
Macierz ma w kolumnie \(j\ge2\)
jedynkę w wierszu \(j-1\),
a pierwsza kolumna jest zerowa.
\end{proof}

\begin{definicja}[Blok Jordana]
\label{def:lin17-jordan-block}
Niech \(\lambda\in K\) oraz \(s\ge1\).
\emph{Blokiem Jordana} rozmiaru \(s\)
odpowiadającym \(\lambda\) nazywamy
macierz \(J_s(\lambda)\in\operatorname{Mat}_{s\times s}(K)\),
której elementy określone są wzorem
\[
 (J_s(\lambda))_{ij}=
 \begin{cases}
 \lambda,&i=j,\\
 1,&j=i+1,\\
 0,&\text{w pozostałych przypadkach}.
 \end{cases}
\]
Dla \(s=1\) mamy \(J_1(\lambda)=(\lambda)\).
Pisząc \(J_s(\lambda)=\lambda I_s+N_s\),
otrzymujemy nilpotentną macierz \(N_s\)
o indeksie nilpotentności \(s\).
\end{definicja}

\begin{przyklad}[Dwa typy bloków nilpotentnych]
\label{prz:lin17-nilpotent-blocks}
Macierze
\[
 N_3=\begin{pmatrix}0&1&0\\0&0&1\\0&0&0\end{pmatrix},
 \qquad
 N_2\oplus N_1=
 \begin{pmatrix}0&1&0\\0&0&0\\0&0&0\end{pmatrix}
\]
mają ten sam wielomian charakterystyczny \(t^3\),
lecz ich wielomiany minimalne są odpowiednio
\(t^3\) i \(t^2\).
Dla \(N_3\) mamy
\(N_3e_3=e_2\), \(N_3e_2=e_1\),
\(N_3e_1=0\); jest to jeden łańcuch długości \(3\).
Drugi operator ma jeden łańcuch długości \(2\)
i jeden łańcuch długości \(1\).
\end{przyklad}

\section{Rozkład nilpotentnego operatora na bloki}
\label{sec:lin17-nilpotent-jordan}

\begin{twierdzenie}[Baza Jordana operatora nilpotentnego]
\label{thm:lin17-nilpotent-jordan}
Niech \(K\) będzie dowolnym ciałem,
\(V\) przestrzenią skończenie wymiarową nad \(K\)
i \(N\in\operatorname{End}_K(V)\)
operatorem nilpotentnym.
Istnieje baza \(V\) złożona z rozłącznych
łańcuchów Jordana operatora \(N\).
W szczególności macierz \(N\) w tej bazie
jest sumą blokową
\[
 [N]\sim J_{s_1}(0)\oplus\cdots\oplus J_{s_r}(0),
 \qquad s_1+\cdots+s_r=\dim V.
\]
Dla \(V=\{0\}\) suma blokowa jest pusta.
\end{twierdzenie}
\begin{proof}
Prowadzimy indukcję względem \(n=\dim V\).
Dla \(n=0\) teza jest oczywista.
Niech \(n\ge1\).
Nilpotentny endomorfizm na niezerowej przestrzeni
nie może być iniektywny.
Istotnie, gdyby \(N\) był iniektywny,
to jego potęga \(N^m\) również byłaby
iniektywna, wbrew \(N^m=0\).
Wobec tego \(W=\operatorname{im}N\)
ma wymiar mniejszy niż \(n\).
Ponieważ \(N(W)\subset W\),
ograniczenie \(N|_W\) jest nilpotentne.
Z założenia indukcyjnego istnieje baza \(W\)
będąca połączeniem łańcuchów
\[
 (w_{i,1},\ldots,w_{i,\ell_i}),
 \qquad 1\le i\le r,
\]
gdzie \(Nw_{i,1}=0\) i
\(Nw_{i,j}=w_{i,j-1}\) dla \(j\ge2\).
Dopuszczamy \(r=0\), jeśli \(W=\{0\}\).

Na \(W\) obraz operatora \(N|_W\)
jest rozpięty przez wektory
\(w_{i,j}\) z \(j<\ell_i\).
Wektory końcowe \(w_{i,\ell_i}\),
dla \(1\le i\le r\), są zatem
liniowo niezależne modulo \(N(W)\).
Dla każdego \(i\) wybierzmy \(v_i\in V\)
tak, że \(Nv_i=w_{i,\ell_i}\);
jest to możliwe, ponieważ \(w_{i,\ell_i}\in W=\operatorname{im}N\).

Pokażemy, że wektory \(v_1,\ldots,v_r\)
są liniowo niezależne modulo \(W\).
Jeśli \(\sum_i a_iv_i\in W\),
to po zastosowaniu \(N\) mamy
\(\sum_i a_iw_{i,\ell_i}\in N(W)\).
Niezależność wektorów końcowych modulo \(N(W)\)
wymusza \(a_i=0\) dla wszystkich \(i\).
Stąd
\[
 U=W\oplus\operatorname{span}\{v_1,\ldots,v_r\},
 \qquad \dim U=\dim W+r.
\]
Podprzestrzeń \(U\) jest niezmiennicza,
bo \(N(W)\subset W\) i \(Nv_i\in W\).
Ponadto \(N(U)=W\):
obrazy \(v_i\) dostarczają końcowych
wektorów łańcuchów w \(W\),
a \(N(W)\) dostarcza wszystkich pozostałych.
Z twierdzenia o rzędzie i jądrze wynika
\[
 \dim(U\cap\ker N)
 =\dim U-\dim N(U)=r.
\]

Uzupełnijmy bazę \(U\cap\ker N\)
do bazy \(\ker N\), dobierając
wektory \(z_1,\ldots,z_q\in\ker N\).
Ponieważ \(\dim\ker N=n-\dim W\),
otrzymujemy \(q=n-\dim W-r\).
Wektory \(z_j\) są liniowo niezależne
modulo \(U\): gdyby ich kombinacja
należała do \(U\), należałaby
do \(U\cap\ker N\), co przeczyłoby
wyborowi uzupełnienia bazy.
Zatem
\[
 V=U\oplus\operatorname{span}\{z_1,\ldots,z_q\},
\]
gdyż prawa strona ma wymiar
\(\dim W+r+q=n\).

Dla każdego \(i\) rodzina
\((w_{i,1},\ldots,w_{i,\ell_i},v_i)\)
stanowi łańcuch Jordana w \(V\),
o jeden wektor dłuższy niż łańcuch w \(W\).
Każdy \((z_j)\) jest łańcuchem długości \(1\).
Połączenie wszystkich tych łańcuchów
jest bazą \(V\), co kończy indukcję.
Jeżeli \(W=0\), mamy \(r=0\) i
\(z_1,\ldots,z_n\) jest dowolną bazą
\(\ker N=V\); rozumowanie obejmuje więc
również operator zerowy.
\end{proof}

\begin{stwierdzenie}[Rozmiary bloków z jąder potęg]
\label{prop:lin17-nilpotent-sizes}
Niech \(N\) będzie nilpotentnym endomorfizmem
skończenie wymiarowej przestrzeni \(V\).
Niech \(d_k=\dim\ker N^k\) dla \(k\ge0\)
oraz \(d_0=0\).
Jeśli w bazie Jordana bloki operatora \(N\)
mają rozmiary \(s_1,\ldots,s_r\), to
\[
 d_k=\sum_{i=1}^r\min(k,s_i),
 \qquad
 d_k-d_{k-1}=\#\{i:s_i\ge k\}
 \quad(k\ge1).
\]
W szczególności liczba bloków o rozmiarze
dokładnie \(k\) wynosi
\[
 b_k=2d_k-d_{k-1}-d_{k+1}.
\]
Wszystkie te liczby są niezależne
od wyboru bazy Jordana.
\end{stwierdzenie}
\begin{proof}
Na bloku \(J_s(0)\) potęga \(N^k\)
zeruje dokładnie pierwszych \(\min(k,s)\)
wektorów standardowej bazy.
Jej jądro ma zatem wymiar \(\min(k,s)\).
Ponieważ jądro potęgi operatora blokowo
diagonalnego jest sumą prostą jąder
odpowiednich bloków, sumowanie daje wzór na \(d_k\).
Różnica \(\min(k,s)-\min(k-1,s)\)
jest równa \(1\) przy \(s\ge k\)
i \(0\) w przeciwnym przypadku.
Otrzymujemy drugi wzór.
Odejmując liczbę bloków rozmiaru co najmniej \(k+1\)
od liczby bloków rozmiaru co najmniej \(k\),
dostajemy trzeci wzór.
Wymiary jąder \(N^k\)
nie zależą od bazy, więc również \(b_k\)
są niezmiennikami operatora.
\end{proof}

\begin{przyklad}[Wykrywanie bloków z wymiarów jąder]
\label{prz:lin17-kernel-sizes}
Załóżmy, że nilpotentny operator \(N\)
na przestrzeni wymiaru \(7\) spełnia
\[
 d_0=0,\qquad d_1=3,\qquad d_2=5,
 \qquad d_3=7,\qquad d_4=7.
\]
Wówczas \(d_1-d_0=3\) oznacza
trzy bloki, \(d_2-d_1=2\)
dwa bloki rozmiaru co najmniej \(2\),
a \(d_3-d_2=2\) dwa bloki
rozmiaru co najmniej \(3\).
Wobec tego dokładnie jeden blok
ma rozmiar \(1\), a pozostałe dwa
rozmiar \(3\); suma rozmiarów wynosi \(1+3+3=7\).
Postać Jordana jest zatem
\(J_3(0)\oplus J_3(0)\oplus J_1(0)\).
\end{przyklad}

\section{Postać Jordana endomorfizmu}
\label{sec:lin17-general-jordan}

\begin{twierdzenie}[Istnienie postaci Jordana]
\label{thm:lin17-jordan}
Niech \(K\) będzie dowolnym ciałem,
\(V\) skończenie wymiarową przestrzenią nad \(K\),
a \(T\in\operatorname{End}_K(V)\).
Jeżeli wielomian minimalny \(m_T\)
rozkłada się nad \(K\) na czynniki liniowe,
to istnieje baza \(V\), w której
macierz operatora \(T\) ma postać
\[
 J=J_{s_1}(\lambda_1)\oplus\cdots
   \oplus J_{s_r}(\lambda_r),
 \qquad \lambda_j\in K,\quad
 s_1+\cdots+s_r=\dim V.
\]
Każdy blok należy do jednej z wartości własnych
operatora \(T\). Równoważnie istnieje
odwracalna macierz \(P\) taka, że
\(P^{-1}[T]P=J\).
Jeśli \(V=\{0\}\), \(J\) jest macierzą pustą.

Warunek rozszczepienia wielomianu minimalnego
jest równoważny rozszczepieniu
wielomianu charakterystycznego \(\chi_T\).
\end{twierdzenie}
\begin{proof}
Z Wniosku~\ref{cor:lin16-generalized}
otrzymujemy rozkład na podprzestrzenie
niezmiennicze
\[
 V=\bigoplus_{\lambda\in\sigma_K(T)}
 E_\lambda^{\mathrm{gen}}(T).
\]
Jeżeli \(V=0\), otrzymujemy pustą sumę bloków.
Niech więc \(V\ne0\).
Na każdej podprzestrzeni
\(W_\lambda=E_\lambda^{\mathrm{gen}}(T)\)
ograniczenie \(T\) jest postaci
\[
 T|_{W_\lambda}=\lambda I_{W_\lambda}+N_\lambda,
 \qquad N_\lambda=(T-\lambda I)|_{W_\lambda}.
\]
Na mocy
Stwierdzenia~\ref{prop:lin16-gen-stabilize}
dla pewnego \(a\ge1\) mamy
\(N_\lambda^a=0\), więc \(N_\lambda\)
jest nilpotentny.
Twierdzenie~\ref{thm:lin17-nilpotent-jordan}
daje bazę \(W_\lambda\) złożoną
z łańcuchów \(N_\lambda\).
W każdym takim łańcuchu
macierz \(T\) jest blokiem \(J_s(\lambda)\).
Łącząc bazy wszystkich \(W_\lambda\),
otrzymujemy bazę \(V\) i żądaną
macierz blokowo diagonalną.

Dla równoważności warunków rozszczepienia
zauważmy, że z twierdzenia
Cayleya--Hamiltona
\(m_T\mid\chi_T\).
Jeśli \(\chi_T\) rozszczepia się na
czynniki liniowe, to rozszczepia się
również każdy jego dzielnik, w szczególności \(m_T\).
W przeciwną stronę, jeśli \(m_T\)
rozszczepia się, to udowodniona część
twierdzenia zapewnia postać Jordana \(J\).
Macierz \(tI-J\) jest blokowo diagonalna,
a każdy blok \(tI-J_s(\lambda)\)
jest górnotrójkątny z przekątną \(t-\lambda\).
Stąd
\(\chi_T(t)=\prod_{j=1}^r(t-\lambda_j)^{s_j}\),
czyli \(\chi_T\) także rozszczepia się.
\end{proof}

\begin{przyklad}[Dwa operatory o tym samym wielomianie charakterystycznym]
\label{prz:lin17-same-char}
Niech \(K=\mathbb C\) i
\[
 A=\begin{pmatrix}2&1&0\\0&2&1\\0&0&2\end{pmatrix},
 \qquad
 B=\begin{pmatrix}2&1&0\\0&2&0\\0&0&2\end{pmatrix}.
\]
Dla obu macierzy
\(\chi(t)=(t-2)^3\).
Niech \(N_A=A-2I\) i \(N_B=B-2I\).
Mamy \(N_A^3=0\), lecz \(N_A^2\ne0\);
natomiast \(N_B^2=0\) i \(N_B\ne0\).
Zatem
\[
 m_A(t)=(t-2)^3,\qquad
 m_B(t)=(t-2)^2.
\]
Pierwsza macierz jest blokiem \(J_3(2)\),
druga sumą \(J_2(2)\oplus J_1(2)\).
Nie są podobne, ponieważ podobieństwo
zachowuje wielomian minimalny.
Ten przykład pokazuje, dlaczego sam
wielomian charakterystyczny nie wystarcza
do wyznaczenia postaci Jordana.
\end{przyklad}

\begin{uwaga}[Kiedy postać Jordana istnieje nad danym ciałem]
\label{rem:lin17-field}
Twierdzenie~\ref{thm:lin17-jordan}
nie wymaga, by \(K\) było algebraicznie domknięte.
Wymaga jedynie, by wielomian charakterystyczny
konkretnego operatora rozszczepiał się nad \(K\).
Nad \(\mathbb C\) warunek jest zawsze spełniony.
Nad \(\mathbb R\) może zawieść:
macierz obrotu
\(\begin{pmatrix}0&-1\\1&0\end{pmatrix}\)
ma wielomian charakterystyczny \(t^2+1\)
i nie jest podobna nad \(\mathbb R\)
do sumy rzeczywistych bloków Jordana.
Zagadnienie postaci kanonicznej
bez rozszczepienia wymaga odrębnego ujęcia.
\end{uwaga}

\section{Jednoznaczność i odczytywanie postaci Jordana}
\label{sec:lin17-uniqueness}

\begin{twierdzenie}[Jednoznaczność rozmiarów bloków Jordana]
\label{thm:lin17-jordan-uniqueness}
Niech \(T\in\operatorname{End}_K(V)\)
ma rozszczepiony wielomian charakterystyczny
i niech \(J\) będzie jego postacią Jordana.
Dla \(\lambda\in\sigma_K(T)\) oraz
\(k\ge0\) oznaczmy
\[
 d_{\lambda,k}
 =\dim\ker((T-\lambda I_V)^k),
 \qquad d_{\lambda,0}=0.
\]
Liczba bloków \(J_s(\lambda)\) rozmiaru
dokładnie \(k\ge1\) jest równa
\[
 b_{\lambda,k}
 =2d_{\lambda,k}
  -d_{\lambda,k-1}
  -d_{\lambda,k+1}.
\]
W szczególności wszystkie rozmiary bloków
i odpowiadające im wartości własne są
jednoznacznie wyznaczone przez \(T\).
Dwie postacie Jordana tego samego operatora
różnią się co najwyżej kolejnością bloków.
\end{twierdzenie}
\begin{proof}
Wybierzmy bazę Jordana \(T\).
Na bloku \(J_s(\lambda)=\lambda I+N_s\)
operator \(T-\lambda I\) działa
jako \(N_s\), więc jądro jego \(k\)-tej potęgi
ma wymiar \(\min(k,s)\).
Dla bloku \(J_s(\mu)\), gdzie \(\mu\ne\lambda\),
operator \(T-\lambda I\) jest odwracalny.
Istotnie, jest równy
\((\mu-\lambda)I+N_s\),
a jego odwrotność wyraża skończony szereg
\[
 ((\mu-\lambda)I+N_s)^{-1}
 =(\mu-\lambda)^{-1}
 \sum_{j=0}^{s-1}
 \bigl(-(\mu-\lambda)^{-1}N_s\bigr)^j.
\]
Po wymnożeniu obie strony
dają identyczność, gdyż \(N_s^s=0\).
Zatem bloki o wartości własnej
różnej od \(\lambda\)
nie wnoszą żadnego wektora
do \(\ker((T-\lambda I)^k)\).
Otrzymujemy więc
\[
 d_{\lambda,k}
 =\sum_{\substack{\text{bloki }J_s(\lambda)}}
          \min(k,s).
\]
Stosując ten sam rachunek różnic
co w Stwierdzeniu~\ref{prop:lin17-nilpotent-sizes},
otrzymujemy wzór na \(b_{\lambda,k}\).
Lewa strona definicji \(d_{\lambda,k}\)
jest niezależna od wyboru bazy,
więc tak samo niezależne są wszystkie
\(b_{\lambda,k}\).
Wartości własne są również niezmiennikami
operatora, co dowodzi jednoznaczności
z dokładnością do kolejności bloków.
\end{proof}

\begin{wniosek}[Wielomian minimalny i diagonalizacja z bloków Jordana]
\label{cor:lin17-jordan-minpoly}
Niech \(T\in\operatorname{End}_K(V)\)
ma rozszczepiony wielomian charakterystyczny.
Jeżeli jego postać Jordana składa się
z bloków \(J_{s_j}(\lambda_j)\), to
\[
 \chi_T(t)=\prod_j(t-\lambda_j)^{s_j},
 \qquad
 m_T(t)=\operatorname{lcm}_j(t-\lambda_j)^{s_j}.
\]
Przy \(V=\{0\}\) pusty iloczyn
i najmniejsza wspólna wielokrotność pustej rodziny
są z definicji równe \(1\).
W szczególności dla każdej wartości własnej
\(\lambda\) wykładnik czynnika
\(t-\lambda\) w \(m_T\)
jest rozmiarem największego bloku
odpowiadającego \(\lambda\).
Operator \(T\) jest diagonalizowalny
nad \(K\) dokładnie wtedy, gdy
wszystkie bloki Jordana mają rozmiar \(1\).
\end{wniosek}
\begin{proof}
Bloki Jordana są górnotrójkątne,
więc wzór na \(\chi_T\)
wynika z iloczynu wyznaczników.
Na bloku \(J_s(\lambda)\)
operator \(J_s(\lambda)-\lambda I_s=N_s\)
ma indeks nilpotentności \(s\).
Wielomian \((t-\lambda)^s\) zeruje blok,
a żadna niższa jego potęga nie zeruje bloku.
Ponieważ wielomian minimalny bloku dzieli
jego wielomian charakterystyczny
\((t-\lambda)^s\) na mocy
Twierdzenia~\ref{thm:lin16-cayley-hamilton},
minimalnym wielomianem tego bloku
musi być \((t-\lambda)^s\).
Wielomian \(p(T)\) zeruje sumę blokową
dokładnie wtedy, gdy zeruje każdy blok,
a więc \(p\) jest podzielny przez
wszystkie \((t-\lambda_j)^{s_j}\).
Monicznym wielomianem dzielącym
każdą taką wspólną wielokrotność
jest ich najmniejsza wspólna wielokrotność,
co daje wzór na \(m_T\).

Jeśli wszystkie bloki mają rozmiar \(1\),
macierz jest diagonalna.
Odwrotnie, jeżeli \(T\) jest diagonalizowalny,
to wielomian minimalny jest
iloczynem różnych czynników liniowych
na mocy
Twierdzenia~\ref{thm:lin16-minpoly-diagonalizable}.
W takim przypadku największy blok
dla każdego \(\lambda\) ma rozmiar \(1\),
zatem wszystkie bloki mają ten rozmiar.
\end{proof}

\begin{przyklad}[Odczytywanie postaci Jordana z ciągów wymiarów]
\label{prz:lin17-jordan-kernels}
Niech \(T\) będzie endomorfizmem
przestrzeni zespolonej wymiaru \(6\),
mającym wartości własne \(1\) i \(3\).
Załóżmy, że
\[
 \begin{gathered}
 d_{1,0}=0,\quad d_{1,1}=2,\quad
 d_{1,2}=3,\quad d_{1,3}=4,\\
 d_{3,0}=0,\quad d_{3,1}=1,\quad
 d_{3,2}=2,\quad d_{3,3}=2.
 \end{gathered}
\]
Dla \(\lambda=1\) różnice
\(d_{1,k}-d_{1,k-1}\)
dla \(k=1,2,3\) wynoszą \(2,1,1\).
Zatem mamy jeden blok długości \(3\)
oraz jeden długości \(1\).
Dla \(\lambda=3\) różnice wynoszą \(1,1,0\):
jest jeden blok długości \(2\).
Ostatecznie
\[
 J=J_3(1)\oplus J_1(1)\oplus J_2(3),
\]
co daje łączny wymiar \(3+1+2=6\).
Wielomiany charakterystyczny i minimalny
są równe odpowiednio
\[
 \chi_T(t)=(t-1)^4(t-3)^2,
 \qquad
 m_T(t)=(t-1)^3(t-3)^2.
\]
\end{przyklad}

\section{Podsumowanie}
\label{sec:lin17-summary}

Nilpotentny endomorfizm przestrzeni
skończenie wymiarowej ma bazę
złożoną z łańcuchów Jordana.
W tej bazie jego macierz jest
sumą blokową nilpotentnych bloków Jordana,
a rozmiary bloków można odczytać
z wymiarów jąder kolejnych potęg.

Jeśli wielomian minimalny operatora \(T\)
rozszczepia się nad \(K\),
rozłożenie \(V\) na uogólnione przestrzenie własne
oraz zastosowanie twierdzenia nilpotentnego
daje postać Jordana \(T\).
Bloki tej postaci są jednoznaczne
z dokładnością do kolejności,
ponieważ ich liczby i rozmiary
można wyznaczyć z wymiarów
\(\ker((T-\lambda I)^k)\).
Wielomian minimalny podaje największy
rozmiar bloku dla danej wartości własnej,
a wielomian charakterystyczny sumę
rozmiarów wszystkich jej bloków.

Postać Jordana uzupełnia diagonalizację:
macierz jest diagonalizowalna
dokładnie wtedy, gdy wszystkie jej
bloki Jordana mają rozmiar jeden.
